% SPDX-FileCopyrightText: 2016-2026 Lari Natri <lari.natri@iki.fi>
% SPDX-License-Identifier: GPL-3.0-or-later
%
% Melody-block parsing, scoped appearance, replay payloads and beat drawing.
% This file is part of the 'ulsbs' package.
%
% Load: ulsbs.sty, after media and before song metadata in the content phase.
% Requires: typography/style, book policy, etoolbox, xparse, expl3 and TikZ.
% Chord/music owners install note rendering and replay integration later.
% State: owns melody policy flags, parser tokens/sequences, boxes, dimensions
% and global replay payload state. Inherits @ as a letter with no enclosing
% group; retains all expl3 spans and the grouped active-hat capture in place.
% Allocations and initial assignments execute now; parsing, dimensions and beat
% drawing execute at use time. Public providecommand appearance defaults may be
% defined before loading or renewed afterward; use the scoped public setters
% for lane/style policy rather than redefining private predicates. Beat display
% permission and switches remain owned by book policy, not this drawing module.

\ProvidesFile{ulsbs-melody.code.tex}[ULSBS melody blocks and beat drawing]

% === Melody blocks and note hints ===

% Melody blocks display transposition-aware encircled note hints in a
% chord definition: \[<C D;E F @ .3em>Am].  Pitches must be uppercase for
% chord transposition mechanism.  `_` advances one empty note slot, `/`
% renders an unpitched circled slash and advances one slot, and `*` marks the
% current slot with a beat.  Negative offsets may protrude left of the shared
% anchor but intentionally never reduce the annotation's rightward advance.
% A leading `!` inside the block renders normally but makes the melody layers
% zero-advance and disables the whitespace minimum for that annotation.
%
% \ulMelodyPrimaryPosition{normal|low},
% \ulMelodySecondaryPosition{above|below|same}, and
% \ulMelodyStyleMapping{normal|swapped} are scoped policies. \ulShowNotes,
% \noteson/\notesoff, \ulShowBeats, and \beatson/\beatsoff remain in effect.

% --- Public melody appearance and layout configuration ---
% Define providecommand defaults before loading or renew afterward. Factors and
% font hooks expand during drawing, not during this declaration phase. Select
% scoped lane/style policies with the public commands below, not private flags.

% Shared font size for melody note circles and their neutral layout slots.
\providecommand{\ulMelodySize}{\scriptsize}

% Setup the style for the melody note name "text". Renew this macro in the
% beginning of your document to set custom style. This will be called each
% time the melody note circle is drawn. This can be set for individual notes,
% too. The base style is \ulMelodySize\mdseries\ulSansHeavyFont; use
% \ulMelodySize for sizing and this drawing-only hook for other appearance.
\providecommand{\ulMelodySymbolStyle}{}

% Setup the style for the small tags put in front of melody note circles. This
% ought to just set a very small font size.
\providecommand{\ulMelodyTagStyle}{\fontsize{3.5pt}{4.5pt}\selectfont}

% Small tag text to be put in front of normal melody note circles. Does not take
% space.
\providecommand{\ulMelodyTagText}{}

% Small tag text to be put in front of alternatively styled melody note circles.
% Does not take space.
\providecommand{\ulMelodyAltTagText}{}

% Factor for the tag text for melody notes is pushed towards the right to be
% closer to the melody note circle. This is relative to 1ex of the melody note
% font, not the font setup in \ulMelodyTagStyle
\providecommand{\ulMelodyTagHorizontalAdjustmentFactor}{.1}

% Factor for the tag text for melody notes is pushed upwards.
% This is relative to 1ex of the melody note font, not the font setup
% in \ulMelodyTagStyle.
\providecommand{\ulMelodyTagVerticalAdjustmentFactor}{.8}

% Factor for the height of the normal melody notes from chord baseline.
% This is relative to 1 ex of the current font: \ulChordFont when called from
% within chord definition as it is meant to.
\providecommand{\ulMelodyNormalHeightFactor}{1.85}

% Factor for the height of the low melody note variant from chord baseline.
% This is relative to 1 ex of the current font: \ulChordFont when called from
% within chord definition as it is meant to.
% Used by the scoped low melody lanes.
\providecommand{\ulMelodyLowHeightFactor}{.3}

% Factor for the height of the high melody note variant from chord baseline.
% This is relative to 1 ex of the current font: \ulChordFont when called from
% within chord definition as it is meant to.
% Used by the high secondary melody lane.
\providecommand{\ulMelodyHighHeightFactor}{3.4}

% Factor for the lane below the low/chord-level melody lane. With the default
% values this continues the same 1.55ex separation used by high/normal/low.
% Selecting primary=low and secondary=below may intentionally overlap lyrics.
\providecommand{\ulMelodyLowestHeightFactor}{-1.25}

% If a visible chord-line annotation is followed by a literal source space,
% reserve at least this much horizontal advance. Wider annotations are
% unchanged. Relative units are evaluated in \ulLyricFontStyle; set to 0pt to disable.
\providecommand{\ulChordWhitespaceMinimumWidth}{.8ex}


% Melody blocks keep musical roles separate from their visual placement and
% styling.  These assignments are deliberately local, like ordinary TeX font
% and color changes, so a verse or explicit group can override them.  `same`
% places secondary on the primary lane; paired voices intentionally overlap.
\newif\ifulsbs@melodyprimarylow\ulsbs@melodyprimarylowfalse
\newif\ifulsbs@melodysecondarybelow\ulsbs@melodysecondarybelowfalse
\newif\ifulsbs@melodysecondarysame\ulsbs@melodysecondarysamefalse
\newif\ifulsbs@melodyswapped\ulsbs@melodyswappedfalse
\newcommand{\ulMelodyPrimaryPosition}[1]{%
  \ifstrequal{#1}{normal}{\ulsbs@melodyprimarylowfalse}{%
    \ifstrequal{#1}{low}{\ulsbs@melodyprimarylowtrue}{%
      \PackageError{ulsbs}{Invalid melody primary position `#1'}%
        {Use \string\ulMelodyPrimaryPosition\{normal\} or \{low\}.}%
    }%
  }%
}
\newcommand{\ulMelodySecondaryPosition}[1]{%
  \ifstrequal{#1}{above}{\ulsbs@melodysecondarybelowfalse\ulsbs@melodysecondarysamefalse}{%
    \ifstrequal{#1}{below}{\ulsbs@melodysecondarybelowtrue\ulsbs@melodysecondarysamefalse}{%
      \ifstrequal{#1}{same}{\ulsbs@melodysecondarybelowfalse\ulsbs@melodysecondarysametrue}{%
        \PackageError{ulsbs}{Invalid melody secondary position `#1'}%
          {Use \string\ulMelodySecondaryPosition\{above\}, \{below\}, or \{same\}.}%
      }%
    }%
  }%
}
\newcommand{\ulMelodyStyleMapping}[1]{%
  \ifstrequal{#1}{normal}{\ulsbs@melodyswappedfalse}{%
    \ifstrequal{#1}{swapped}{\ulsbs@melodyswappedtrue}{%
      \PackageError{ulsbs}{Invalid melody style mapping `#1'}%
        {Use \string\ulMelodyStyleMapping\{normal\} or \{swapped\}.}%
    }%
  }%
}

% Prints the pretag for melody note circle, if one exists
\newcommand{\ulsbs@mntag}{%
  \begingroup%
    \hspace{\dimexpr\ulMelodyTagHorizontalAdjustmentFactor ex\relax}%
    \raisebox{\dimexpr\ulMelodyTagVerticalAdjustmentFactor ex\relax}{%
        \llap{\ulMelodyTagStyle\ifulsbs@altmnstyle\ulMelodyAltTagText\else\ulMelodyTagText\fi}%
    }%
  \endgroup%
}

% Base melody-note symbol style, followed by \ulMelodySymbolStyle:
\newcommand{\ulsbs@mnsymbolstyle}{%
  \ulMelodySize\mdseries\ulSansHeavyFont%
  \ifulsbs@altmnstyle%
    \color{ulMelodyNoteAltColor}%
  \else%
    \color{ulMelodyNoteColor}%
  \fi%
  \ulMelodySymbolStyle%
}

% Sharp and flat symbols for note circles: adjust their height and kerning
% for a better fit within the circle.
\newcommand{\ulsbs@mnsharpsymbol}{\raise-.25ex\hbox{\textsuperscript{{\kern-0.05em\ulSymbolFont♯}}}}%


% Outputs an encircled melody note, internal
%   - optional parameter: rise the note above text (chord line) by this amount,
%                         defaults to \ulMelodyNormalHeightFactor ex
%   - mandatory parameter: the note, must be in UPPER CASE for transposing to
%     work
%   * Example usage: \ulsbs@note[0.5em]{C#}
%   * To take no horizontal space, call this macro within \makebox[0pt][l]{}
\NewDocumentCommand{\ulsbs@note}{O{\dimexpr\ulMelodyNormalHeightFactor ex\relax} m}{}

% --- Melody-block parser state ---
%
% The syntax is parsed after the chord scanner collects one \[...] argument.
% Angle brackets retain their ordinary catcode and have no meaning elsewhere.
% Parsed primary and secondary items are stored as expl3 sequences so renderer
% layout is independent of the textual spelling (compact or space-separated).
% The one-token `!` modifier is dispatched separately below to avoid expanding
% active accidental tokens before they reach this content parser.
\newsavebox{\ulsbs@melodyprimarybox}
\newsavebox{\ulsbs@melodysecondarybox}
\newsavebox{\ulsbs@melodybeatbox}
\newsavebox{\ulsbs@melodylayerbox}
\newdimen\ulsbs@melodyslotwidth
\newdimen\ulsbs@melodyoffset
\newdimen\ulsbs@melodybeatwidth
\newdimen\ulsbs@melodybeatextent
\newdimen\ulsbs@melodycurrentextent
\newif\ifulsbs@annotationvisible
\newif\ifulsbs@annotationfullclearance
\newif\ifulsbs@annotationminimumdisabled
\newif\ifulsbs@melodyzeroadvance

\ExplSyntaxOn
\tl_new:N \l__ulsbs_melody_primary_tl
\tl_new:N \l__ulsbs_melody_secondary_tl
\tl_new:N \l__ulsbs_melody_offset_tl
\tl_new:N \l__ulsbs_melody_at_tl
\seq_new:N \l__ulsbs_melody_parts_seq
\seq_new:N \l__ulsbs_melody_primary_seq
\seq_new:N \l__ulsbs_melody_secondary_seq
\seq_new:N \l__ulsbs_melody_lex_seq
\cs_generate_variant:Nn \seq_set_split:Nnn { NVn, NVV }
\cs_generate_variant:Nn \str_case:nnF { en }
\cs_generate_variant:Nn \regex_match:nnT { nV }
% Report parser errors with the complete original block for useful diagnostics.
\cs_new_protected:Npn \__ulsbs_melody_error:n #1
  { \PackageError{ulsbs}{Malformed~melody~block~`\tl_use:N\l__ulsbs_melody_raw_tl'}{#1} }
\tl_new:N \l__ulsbs_melody_raw_tl
\tl_new:N \l__ulsbs_melody_content_tl
% Parse separators first: one optional @ dimension and one optional ; voice
% separator. The remaining voice fields are then tokenized independently.
\cs_new_protected:Npn \__ulsbs_melody_parse:n #1
 {
  \tl_set:Nn \l__ulsbs_melody_raw_tl {#1}
  \tl_set:Nn \l__ulsbs_melody_content_tl {#1}
  \tl_trim_spaces:N \l__ulsbs_melody_content_tl
  \regex_match:nVT { @ \s* \Z } \l__ulsbs_melody_content_tl
   {\__ulsbs_melody_error:n {The~@~separator~must~be~followed~by~a~TeX~dimension.}}
  \regex_match:nVT { @ .* @ } \l__ulsbs_melody_content_tl
   {\__ulsbs_melody_error:n {A~melody~block~may~contain~only~one~@~offset.}}
  \tl_set:Nx \l__ulsbs_melody_at_tl {\char_generate:nn {64} {12}}
  \seq_set_split:NVV \l__ulsbs_melody_parts_seq \l__ulsbs_melody_at_tl \l__ulsbs_melody_content_tl
  \int_compare:nNnTF {\seq_count:N \l__ulsbs_melody_parts_seq} > {2}
   {\__ulsbs_melody_error:n {A~melody~block~may~contain~only~one~@~offset.}}
   {
    \seq_pop_left:NN \l__ulsbs_melody_parts_seq \l__ulsbs_melody_primary_tl
    \tl_trim_spaces:N \l__ulsbs_melody_primary_tl
    \tl_clear:N \l__ulsbs_melody_offset_tl
    \seq_if_empty:NTF \l__ulsbs_melody_parts_seq
     {\tl_if_in:VnT \l__ulsbs_melody_raw_tl {@}
       {\__ulsbs_melody_error:n {The~@~separator~must~be~followed~by~a~TeX~dimension.}}}
     { \seq_pop_left:NN \l__ulsbs_melody_parts_seq \l__ulsbs_melody_offset_tl
       \tl_trim_spaces:N \l__ulsbs_melody_offset_tl
       \tl_if_empty:NT \l__ulsbs_melody_offset_tl
        {\__ulsbs_melody_error:n {The~@~separator~must~be~followed~by~a~TeX~dimension.}} }
    \seq_set_split:NnV \l__ulsbs_melody_parts_seq {;} \l__ulsbs_melody_primary_tl
    \int_compare:nNnTF {\seq_count:N \l__ulsbs_melody_parts_seq} > {2}
     {\__ulsbs_melody_error:n {A~melody~block~may~contain~only~one~;~separator.}}
     { \seq_pop_left:NN \l__ulsbs_melody_parts_seq \l__ulsbs_melody_primary_tl
       \tl_trim_spaces:N \l__ulsbs_melody_primary_tl
       \tl_clear:N \l__ulsbs_melody_secondary_tl
       \seq_if_empty:NF \l__ulsbs_melody_parts_seq
        {\seq_pop_left:NN \l__ulsbs_melody_parts_seq \l__ulsbs_melody_secondary_tl
         \tl_trim_spaces:N \l__ulsbs_melody_secondary_tl}
       \seq_clear:N \l__ulsbs_melody_primary_seq
       \seq_clear:N \l__ulsbs_melody_secondary_seq
       \__ulsbs_melody_lex:NV \l__ulsbs_melody_primary_seq \l__ulsbs_melody_primary_tl
       \__ulsbs_melody_lex:NV \l__ulsbs_melody_secondary_seq \l__ulsbs_melody_secondary_tl
        % Keep dimensions unexpanded until the renderer knows there is melody ink.
        \ulsbs@melodyoffset=0pt\relax
     }
   }
 }
% The lexer ignores spaces and recognizes pitches one character at a time.
% Accidentals are appended to the pending pitch without stringifying their
% active tokens, which keeps chord transposition intact.
\tl_new:N \l__ulsbs_melody_pending_tl
\cs_new_protected:Npn \__ulsbs_melody_lex_commit:
 {\tl_if_empty:NF \l__ulsbs_melody_pending_tl
   {\seq_put_right:NV \l__ulsbs_melody_lex_seq \l__ulsbs_melody_pending_tl\tl_clear:N \l__ulsbs_melody_pending_tl}}
\cs_new_protected:Npn \__ulsbs_melody_lex_token:n #1
 {\str_if_eq:eeTF {\tl_to_str:n {#1}} {\c_space_tl} {}
   {\str_if_eq:eeTF {\tl_to_str:n {#1}} {\c_hash_str}
      {\tl_if_empty:NTF \l__ulsbs_melody_pending_tl
        {\__ulsbs_melody_error:n {An~accidental~must~immediately~follow~a~pitch.}}
        {\tl_put_right:Nn \l__ulsbs_melody_pending_tl {#1}\__ulsbs_melody_lex_commit:}}
      {\str_if_eq:eeTF {\tl_to_str:n {#1}} {&}
        {\tl_if_empty:NTF \l__ulsbs_melody_pending_tl
          {\__ulsbs_melody_error:n {An~accidental~must~immediately~follow~a~pitch.}}
          {\tl_put_right:Nn \l__ulsbs_melody_pending_tl {#1}\__ulsbs_melody_lex_commit:}}
        {\__ulsbs_melody_lex_commit:
         \str_case:enF {\tl_to_str:n {#1}}
          {{A}{\tl_set:Nn \l__ulsbs_melody_pending_tl {#1}}{B}{\tl_set:Nn \l__ulsbs_melody_pending_tl {#1}}{C}{\tl_set:Nn \l__ulsbs_melody_pending_tl {#1}}{D}{\tl_set:Nn \l__ulsbs_melody_pending_tl {#1}}{E}{\tl_set:Nn \l__ulsbs_melody_pending_tl {#1}}{F}{\tl_set:Nn \l__ulsbs_melody_pending_tl {#1}}{G}{\tl_set:Nn \l__ulsbs_melody_pending_tl {#1}}{_}{\seq_put_right:Nn \l__ulsbs_melody_lex_seq {#1}}{*}{\seq_put_right:Nn \l__ulsbs_melody_lex_seq {#1}}{/}{\seq_put_right:Nn \l__ulsbs_melody_lex_seq {#1}}}
          {\__ulsbs_melody_error:n {Each~item~must~be~A--G~with~optional~sharp/flat,~_,~*,~or~/.}}}}}}
\cs_new_protected:Npn \__ulsbs_melody_lex:Nn #1#2
 {\seq_clear:N \l__ulsbs_melody_lex_seq\tl_clear:N \l__ulsbs_melody_pending_tl
  \tl_map_function:nN {#2} \__ulsbs_melody_lex_token:n
  \__ulsbs_melody_lex_commit:\seq_set_eq:NN #1 \l__ulsbs_melody_lex_seq}
\cs_generate_variant:Nn \__ulsbs_melody_lex:Nn { NV }
\cs_new_eq:NN \ulsbs@melodyparsecontent \__ulsbs_melody_parse:n
\ExplSyntaxOff

% Dispatch the optional leading ! without expanding the remaining token list.
% Keeping this one-token modifier parser outside expl3 preserves active
% accidental tokens verbatim for the content lexer and transposition.
\def\ulsbs@melodyparse#1{%
  \expandafter\ulsbs@melodyparsemodifier#1\ulsbs@melodyend%
}
\def\ulsbs@melodyparsemodifier{%
  \@ifnextchar!\ulsbs@melodyparsezero\ulsbs@melodyparsenormal%
}
\def\ulsbs@melodyparsezero!#1\ulsbs@melodyend{%
  \ulsbs@melodyzeroadvancetrue
  \ulsbs@annotationminimumdisabledfalse
  \ulsbs@melodyparsecontent{#1}%
}
\def\ulsbs@melodyparsenormal#1\ulsbs@melodyend{%
  \ulsbs@melodyzeroadvancefalse
  \ulsbs@annotationminimumdisabledfalse
  \ulsbs@melodyparsecontent{#1}%
}

% Resolve the semantic primary/secondary role to the currently selected visual
% style immediately before drawing each note.
\newcommand{\ulsbs@melodystyle}[1]{%
  \ifulsbs@melodyswapped%
    \ifstrequal{#1}{primary}{\ulsbs@altmnstyletrue}{\ulsbs@altmnstylefalse}%
  \else%
    \ifstrequal{#1}{primary}{\ulsbs@altmnstylefalse}{\ulsbs@altmnstyletrue}%
  \fi%
}
\ExplSyntaxOn
% Beat drawing is visibility-aware. The legacy .30ex lead is visual only: it
% shifts the glyph but is deliberately excluded from reserved layer width.
\newcommand{\ulsbs@melodybeatdraw}{%
  \ifulShowBeats%
    {\ulLyricFontStyle\kern\ulsbs@bmleadspace}\ulsbs@bmdraw%
  \fi%
}
\cs_new_protected:Npn \__ulsbs_melody_sequence:Nnn #1#2#3
 { \seq_map_inline:Nn #1
   { \str_if_eq:eeF {\tl_to_str:n {##1}} {*} { \str_if_eq:eeTF {\tl_to_str:n {##1}} {_}
       {\makebox[\ulsbs@melodyslotwidth][c]{}}
       {\makebox[\ulsbs@melodyslotwidth][c]{#2\ulsbs@note[#3]{##1}}} } } }
\cs_new_protected:Npn \__ulsbs_melody_beats:N #1
 { \seq_map_inline:Nn #1
   { \str_if_eq:eeTF {\tl_to_str:n {##1}} {*} {\hbox~to~0pt{\ulsbs@melodybeatdraw\hss}}
       {\hskip\ulsbs@melodyslotwidth} } }
\cs_new_protected:Npn \__ulsbs_melody_beat_extent:N #1
 {
  \ulsbs@melodycurrentextent=0pt\relax
  \seq_map_inline:Nn #1
   { \str_if_eq:eeTF {\tl_to_str:n {##1}} {*}
      {\@tempdima=\dimexpr\ulsbs@melodycurrentextent+\ulsbs@melodybeatwidth\relax
       \ifdim\@tempdima>\ulsbs@melodybeatextent\ulsbs@melodybeatextent=\@tempdima\fi}
      {\advance\ulsbs@melodycurrentextent by \ulsbs@melodyslotwidth\relax} }
 }
\cs_new_protected:Npn \__ulsbs_melody_mark_note_visible:N #1
 {
  \seq_map_inline:Nn #1
   {
    \str_if_eq:eeF {\tl_to_str:n {##1}} {*}
     {
      \str_if_eq:eeF {\tl_to_str:n {##1}} {_}
       {\ulsbs@annotationvisibletrue\seq_map_break:}
     }
   }
 }
\ExplSyntaxOff

% --- Integration with chord memorization and replay ---
%
% A zero-output payload is stored in the chord register beside ordinary
% chord tokens. On replay it recovers the raw melody block, which is parsed and
% rendered again under the replay location's visibility, placement and style.
\newif\ifulsbs@melodypayloadempty
\newif\ifulsbs@melodyhadpayload
\global\let\ulsbs@melodyreplayraw\@empty
% #1 = raw block; #2 = 1 when this annotation has no chord remainder.
\def\ulsbs@melodyreplaypayload#1#2{%
  \global\ulsbs@melodyhadpayloadtrue
  \global\def\ulsbs@melodyreplayraw{#1}%
  \ifnum#2=1 \global\ulsbs@melodypayloademptytrue\else\global\ulsbs@melodypayloademptyfalse\fi%
}

\begingroup\catcode`\^=\active\global\let\ulsbs@melodyactivehat=^\endgroup

% === Beat marks ===

% Beat marks are written as `*` in a melody block, for example \[<* C>Am].
% They are small dots just below the chord baseline and share the block's
% horizontal offset and annotation width calculation.
%
% To disable showing of the beats for the whole document, set \ulShowBeats{false}
% in the main document preamble (after including this package).
%
% Call \beatsoff command between verses to disable showing of beat marks in the
% following verses of the song. \beatson puts them back for the next verses
% (unless \ulShowBeats{false} is set in the main doc).
%

% Private drawing defaults: absolute dimensions except the font-relative lead.
\providecommand*{\ulsbs@bmradius}{1.1pt}    % mark radius
\providecommand*{\ulsbs@bmdrop}{1.1pt}      % how far mark's bottom goes below baseline
\providecommand*{\ulsbs@bmdepth}{.7pt}      % mark's reported depth
\providecommand*{\ulsbs@bmleadspace}{.30ex} % leading space before mark (rel. to \ulLyricFontStyle)

% Draws the beat mark.
%
% NOTE: without raisebox, the TikZ image's bottom would be placed on baseline,
% and left would be put horizontally at current position.
% With raisebox{-1.1pt} the bottom of the circle goes 1.1pt below baseline.
% The default ulChordLineParams has lineskiplimit=.5pt and lineskip=.5pt, which
% ensures there is always a minimum of .5pt space between chord line's bottom
% and lyric line's top, when we set the beat mark to have .7pt depth (when
% it really is 1.1pt), so it goes .4 pt deeper than informed, so there is
% always at least .1pt between lyric top and beat mark bottom.
% TODO: use saved box (had problems with dynamic color when tried)
\newcommand*{\ulsbs@bmdraw}{%
  \begingroup%
  \color{ulMeasureBarColor}%
  \raisebox{-\ulsbs@bmdrop}[0pt][\ulsbs@bmdepth]{%
    \tikz\fill (0,0) circle [radius=\ulsbs@bmradius];%
  }%
  \endgroup%
}

% Beat permission, grouped requests and short switches belong to ulsbs-visibility.code.tex.
